\GalSourcePageBoundary{086}{L0085}{57}{57}
des valeurs de $k_{1}$~et~$k_{2}$ finies et déterminées. Ainsi la condition
pour qu'une pareille transformation soit réellement une substitution
est que $m_{1}n_{2} - m_{2}n_{1}$ ne soit ni nul ni divisible par le module~$p$,
ce qui est la même chose.

Je dis maintenant que, bien que ce groupe à substitutions
linéaires n'appartienne pas toujours, comme on le verra, à des
équations solubles par radicaux, il jouira toutefois de cette propriété,
que, si dans une quelconque de ses substitutions il y a $n$~lettres
de fixes, $n$~divisera le nombre des lettres. Et, en effet, quel
que soit le nombre des lettres qui restent fixes, on pourra exprimer
cette circonstance par des équations linéaires qui donneront
tous les indices de l'une des lettres fixes, au moyen d'un certain
nombre d'entre eux. Donnant à chacun de ces indices, restés arbitraires,
$p$~valeurs, on aura $p^{m}$~systèmes de valeurs, $m$~étant un certain
nombre. \GalSourceError{Dans le cas qui nous occupe, $m$~est nécessairement
$< 2$ et se trouve par conséquent être $0$ ou~$1$.}{W10-SE001} Donc le nombre des
substitutions ne saurait être plus grand que
\[
p^{2}(p^{2} - 1)(p^{2} - p).
\]

Ne considérons maintenant que les substitutions linéaires où la
lettre~$a_{0,0}$ ne varie pas; si, dans ce cas, nous trouvons le nombre
total des permutations du groupe qui contient toutes les substitutions
linéaires possibles, il nous suffira de multiplier ce nombre
par~$p^{2}$.

Or, premièrement, en substituant $p$ à l'indice~$k_{2}$, toutes les
substitutions de la forme
\[
(a_{k_{1}, k_{2}},\ a_{m_{1}k_{1}, k_{2}})
\]
donneront en tout $p - 1$ substitutions. On en aura $p^{2} - p$ en
ajoutant au terme~$k_{2}$ le terme~$m_{2}k_{1}$, ainsi qu'il suit:
\[
\tag{m'}
(k_{1},\ k_{2},\ m_{1}k_{1},\ m_{2}k_{1} + k_{2}).
\]

D'un autre côté, il est aisé de trouver un groupe linéaire
de $p^{2} - 1$ permutations, tel que, dans chacune de ses substitutions,
toutes les lettres, à l'exception de~$a_{0,0}$, varient. Car, en
remplaçant le double indice $k_{1}$,~$k_{2}$ par l'indice simple $k_{1} + ik_{2}$,
